\section{Preliminaries}\label{sec:prelims}

We write $[k]=\{1,\dots,k\}$, $\log$ for the natural logarithm, and $\Hscale:=\loglog n$.
For an integer $a\ge 0$ and a real $y$ we use the falling factorial $(y)_a:=y(y-1)\cdots(y-a+1)$.
Throughout, ``\whp'' means with probability $1-o(1)$ as $n\to\infty$, and constants hidden in
$\Oh(\cdot),\Om(\cdot),\Th(\cdot)$ are absolute unless a dependence is shown.

\subsection{The model and the two discrepancy quantities}\label{sec:model}

We work in the average-case online vector balancing model of Altschuler and Tikhomirov~\cite{AT2025},
extended from two colors to $k$ colors.

\begin{definition}[The sparse i.i.d.\ model $\mu_{d,n}$]\label{def:model}
Fix integers $n,d,k$ and $T=\Th(n)$, with $2\le d\le \Hscale^2/\log\Hscale$ and $2\le k$.
The arrivals $v_1,\dots,v_T\in\{0,1\}^n$ are drawn independently and uniformly from the set of $0/1$ vectors
with exactly $d$ ones; equivalently, the \emph{support} $\supp(v_t)\subseteq[n]$ is a uniformly random
$d$-subset, independent across $t$.
\end{definition}

The sparsity bound $d\le \Hscale^2/\log\Hscale=(\loglog n)^2/\log\log\log n$ is the regime in which the two-color
value is independent of $d$~\cite{AT2025}; we keep it throughout.
A coloring assigns each arrival a color $\chi(t)\in[k]$. In the \emph{online} model the assignment is irrevocable
and made knowing only $v_1,\dots,v_t$; in the \emph{offline} model the coloring may depend on all of $v_1,\dots,v_T$.
For color $c\in[k]$ let $\load_c(t)=\sum_{s\le t:\ \chi(s)=c}v_s\in\Z_{\ge0}^n$ be its load. For coordinate $i\in[n]$
the \emph{profile} is $x_i(t)=(\load_{1,i}(t),\dots,\load_{k,i}(t))$, and the \emph{pairwise range} is
\[
  \rangeop_i(t)\ :=\ \max_{c\in[k]}\load_{c,i}(t)\ -\ \min_{c\in[k]}\load_{c,i}(t).
\]
We use two discrepancy quantities, following the online/offline split of~\cite{AT2025}. The online quantity is the
\emph{prefix} discrepancy: it must be small at every prefix, since the coloring is irrevocable. The offline quantity is the
\emph{final} discrepancy of the resulting set system, the standard offline notion. We study their optimal values:
\[
  \ondisc_k(\mu_{d,n})\ :=\ \min_{\text{online }\chi}\ \max_{t\le T,\,i}\ \rangeop_i(t),
  \qquad
  \offdisc_k(\mu_{d,n})\ :=\ \min_{\text{offline }\chi}\ \max_{i}\ \rangeop_i(T),
\]
each understood \whp\ (and, where stated, in expectation). We abbreviate $\ondisc_k=\ondisc_k(\mu_{d,n})$ and
$\offdisc_k=\offdisc_k(\mu_{d,n})$. The online prefix discrepancy dominates the online final discrepancy, which in turn
dominates the offline final discrepancy, so $\ondisc_k\ge\offdisc_k$ always.

\subsection{Scales and the threshold function}\label{sec:scales}

Two scales drive the paper. The \emph{occupancy}
\[
  \lam\ :=\ \frac{Td}{nk}\ =\ \Th\!\left(\frac{d}{k}\right)
\]
is the expected number of arrivals of a fixed color that cover a fixed coordinate, and
$\Lscale:=\log(ek)$. The offline value is governed by the threshold function
\begin{equation}\label{eq:Psi}
  \Psifn(\lam,\Lscale)\ :=\ 1\ +\ \sqrt{\lam\Lscale}\ +\ \frac{\Lscale}{\log\!\big(e+\Lscale/\lam\big)}.
\end{equation}
The three summands separate three regimes of $\Psifn(\lam,\log k)$, corresponding to the Gaussian, large-deviation,
and bounded-occupancy behaviour of $k$ independent Poisson loads (see \Cref{sec:offline-lower}):
\begin{equation}\label{eq:Psi-regimes}
  \Psifn(\lam,\log k)\ \asymp\
  \begin{cases}
    \sqrt{\lam\log k}, & \lam\gtrsim\log k,\\[2pt]
    \dfrac{\log k}{\log\big((\log k)/\lam\big)}, & 1\ll\lam\ll\log k,\\[6pt]
    \dfrac{\log k}{\loglog k}, & \lam=\Th(1).
  \end{cases}
\end{equation}
It is worth reading off the three cases one at a time, because each one is the answer to the same question asked at a
different density. The question is: by how much does the largest of $k$ independent $\Pois(\lam)$ loads exceed the typical
load $\lam$? In the dense regime $\lam\gtrsim\log k$ each load is already concentrated to within $\Th(\sqrt\lam)$ of its
mean, so the maximum over $k$ of them sits about $\sqrt{\log k}$ standard deviations out, and the spread is the Gaussian
value $\sqrt{\lam\log k}$. As $\lam$ drops below $\log k$ the Poisson tail stops looking Gaussian: the relevant deviation
$t$ is now large compared with the mean, the tail decays like $e^{-t\log(t/\lam)}$ instead of $e^{-t^2/\lam}$, and solving
for the $t$ at which the expected number of colors past $\lam+t$ falls to one gives the Bennett value
$(\log k)/\log\big((\log k)/\lam\big)$. Pushing $\lam$ all the way to $\Th(1)$ freezes this at its extreme: the spread is
$(\log k)/\loglog k$, the largest occupancy of $k$ balls thrown into one bin, which no longer depends on $\lam$ at all.
The two non-constant summands of \eqref{eq:Psi} are these two answers, and they are written so that whichever regime one is
in, the larger summand is the one that governs.

The seam between the Gaussian and Bennett pieces sits at $\lam\asymp\Lscale$. There both $\sqrt{\lam\Lscale}$ and
$\Lscale/\log(e+\Lscale/\lam)$ are $\Th(\Lscale)$, so the two summands meet at the same order and neither has a corner: the
$\log(e+\cdot)$ in the denominator is exactly the device that lets the second summand decay smoothly from $\Th(\Lscale)$ at
the seam up toward the bounded-occupancy plateau as $\lam$ shrinks. Consequently $\Psifn$ is continuous in $(\lam,\Lscale)$,
the maximum of its two parts tracks the true Poisson spread across all three regimes, and the leading $1$ keeps $\Psifn\ge1$
even when both deviations are tiny.

We will also use a per-coordinate \emph{rank} that records both the range and how many colors attain the maximum. The range
alone is a blunt instrument for a recursion: two profiles can have the same range while one is a lone spike and the other
has many colors tied at the top, and these behave differently when split further. The rank below folds the integer range
together with the multiplicity of the maximum into a single monotone quantity, so that the dominant term $(k-1)(\rangeop-1)$
recovers the range up to the additive correction $\maxmult\in[1,k]$ and the displayed inequality reads the range back out.
For a profile $x\in\Z_{\ge0}^k$ with range $\rangeop(x)$ and maximum multiplicity $\maxmult(x):=|\arg\max_c x_c|$, set
\begin{equation}\label{eq:rank}
  \rank(x)\ :=\
  \begin{cases}
    0, & \rangeop(x)=0,\\
    (k-1)(\rangeop(x)-1)+\maxmult(x), & \rangeop(x)\ge 1,
  \end{cases}
  \qquad\text{so that}\qquad
  \rangeop(x)\ \le\ 1+\frac{\rank(x)}{k-1}.
\end{equation}

\subsection{Negative association and strong Rayleigh measures}\label{sec:NA}

The offline upper bound rests on a single draw of the random incidence matrix being well behaved for every column
subset simultaneously. We want to control sums of many incidence entries and to take a union over exponentially many
events on \emph{one} sample, which means we cannot afford the positive correlations that an independence assumption would
have ruled out for free. The columns of $\mu_{d,n}$ are independent, but within a column the $n$ coordinates are not: fixing
$d$ ones makes the entries compete, so a one in one coordinate makes a one elsewhere slightly less likely. This competition
is the helpful direction, and negative association is the tool that makes it usable.

\begin{definition}[Negative association~\cite{JoagDevProschan1983}]\label{def:NA}
Random variables $(Y_1,\dots,Y_m)$ are \emph{negatively associated} (NA) if for every two disjoint index sets
$I,J\subseteq[m]$ and every pair of coordinatewise-nondecreasing functions $f,g$,
\[
  \E\!\big[f(Y_I)\,g(Y_J)\big]\ \le\ \E\big[f(Y_I)\big]\,\E\big[g(Y_J)\big].
\]
\end{definition}

We use three standard facts, and they chain together exactly to fit the structure of the incidence matrix. First, the
indicator vector of a uniformly random $d$-subset of $[n]$ is strongly Rayleigh, hence NA across the $n$
coordinates~\cite{BorceaBrandenLiggett2009}. Strong Rayleigh is the strongest of the common negative-dependence notions,
and the point of invoking it is that it is robust: it is preserved under the operations we are about to perform, whereas a
bare NA hypothesis can be fragile. Second, a vector formed by concatenating independent NA vectors is NA, which is what lets
us pass from a single column to the whole matrix, since the columns are drawn independently. Third, if
$(Z_1,\dots,Z_m)$ is NA and $g_1,\dots,g_r$ are coordinatewise-nondecreasing functions acting on \emph{disjoint} blocks of
the coordinates, then $(g_1(Z),\dots,g_r(Z))$ is NA; applied with $g_i$ the row-count over a fixed column subset, this hands
us the joint negative association of the quantities the structural lemma actually controls. The single consequence we need
out of all this is that NA grants the upper-tail moment-generating-function factorization
$\E\big[\prod_i e^{tY_i}\big]\le\prod_i\E\big[e^{tY_i}\big]$ for $t>0$ and nondecreasing $Y_i$, exactly as for
independent variables. That one inequality is what we use: it is the only place independence would otherwise have entered,
and it is precisely the step that turns a per-coordinate tail bound into a Chernoff bound for the sum, so that the union
over column subsets goes through on a single draw. We never need a matching lower-tail or two-sided statement, only this
upper-tail factorization. The relevant closure properties are recalled in \Cref{app:prelims}.

\subsection{Partial coloring and Poisson tails}\label{sec:tools}

The offline construction halves color classes with the partial-coloring lemma of Lovett and Meka.

\begin{lemma}[Lovett--Meka partial coloring~\cite{LovettMeka2015}]\label{lem:LM}
Let $x_0\in[-1,1]^m$ and let $v_1,\dots,v_r\in\R^m$ have thresholds $c_1,\dots,c_r\ge 0$ with
$\sum_{j} e^{-c_j^2/16}\le m/16$. Then there is $x\in[-1,1]^m$ with at least $m/2$ coordinates in $\{-1,1\}$ and
$|\langle v_j,\,x-x_0\rangle|\le c_j\|v_j\|_2$ for every $j$. A threshold $c_j=0$ enforces the linear equality
$\langle v_j,\,x-x_0\rangle=0$ exactly.
\end{lemma}

The Lovett--Meka lemma is the workhorse of the upper bound: it produces a partial coloring meeting many linear constraints
at once, and the threshold $c_j=0$ option is what we use to pin part sizes exactly while the nonzero thresholds keep the row
counts under control. We invoke it repeatedly, peeling off half the alive columns each time.

The offline lower bound and the broadening analysis use the Bennett--Poisson tail. This is the quantitative form of the two
regimes already seen in $\Psifn$: the small-deviation branch $t^2/\mu$ is the Gaussian rate that drives the
$\sqrt{\lam\Lscale}$ summand, and the large-deviation branch $t\log(1+t/\mu)$ is the Bennett rate that drives the second
summand. For $X\sim\Pois(\mu)$ and $t\ge0$,
\begin{equation}\label{eq:poisson-tail}
  -\log\Pr\big(X\ge \mu+t\big)\ \asymp\
  \begin{cases}
    t^2/\mu, & t\lesssim\mu,\\
    t\,\log(1+t/\mu), & t\gtrsim\mu.
  \end{cases}
\end{equation}
Equating the rate in \eqref{eq:poisson-tail} to $\Lscale$ and solving for $t$ yields
$t\asymp\sqrt{\lam\Lscale}+\Lscale/\log(e+\Lscale/\lam)$, the two non-constant summands of $\Psifn$; this is the
calculation behind \eqref{eq:Psi-regimes} (carried out in \Cref{sec:offline-lower}).
